
ce qui
montre que @@M0@@.

\textbf{Corollaire 6.12.} \uline{Si} @@M1@@ \uline{est de caractéristique nulle}, \uline{et} @@M2@@ \uline{une} @@M3@@-@@M4@@-\uline{algèbre} \uline{de filtration discrète}, \uline{alors la classe de Chern complétée} :
@@M5@@
\uline{est un isomorphisme}.

La démonstration est contenue dans celle de 6.11.

\textbf{Définition 6.13.} \uline{Soit} @@M6@@ \uline{une} @@M7@@-\uline{algèbre graduée}, \uline{et} @@M8@@ \uline{une} @@M9@@-@@M10@@-\uline{algèbre} \uline{augmentée}~; \uline{on appelle théorie de Chern de} @@M11@@ \uline{à valeurs dans} @@M12@@ \uline{tout} @@M13@@-\uline{homomorphisme} : @@M14@@, \uline{compatible avec les augmentations}.

On définit un foncteur sur la catégorie des @@M15@@-algèbres graduées, à valeurs dans la catégorie des ensembles, en associant à toute @@M16@@-algèbre graduée @@M17@@ l'ensemble des théories de Chern de @@M18@@ à valeurs dans @@M19@@, soit @@M20@@.

\textbf{Proposition 6.14.} \uline{Le foncteur défini en} 6.13 \uline{est représentable}. \uline{Si} @@M21@@ \uline{est} \uline{de caractéristique nulle}, \uline{et} @@M22@@ \uline{de filtration discrète}, \uline{il est représenté} \uline{par le gradué associé à} @@M23@@.

Montrons d'abord le premier point. On considère la @@M24@@-algèbre graduée engendrée par des générateurs @@M25@@, avec @@M26@@, @@M27@@, @@M28@@ étant de degré @@M29@@, et on fait son quotient par les relations suivantes :
@@M30@@
